% !TEX program = pdflatex
\documentclass[aspectratio=169,10pt]{beamer}

% =========================
% PAQUETES
% =========================
\usepackage[spanish,es-noshorthands]{babel}
\usepackage[T1]{fontenc}
\usepackage[utf8]{inputenc}
\usepackage{lmodern}
\usepackage{amsmath,amssymb,amsfonts}
\usepackage{graphicx}
\usepackage{booktabs}
\usepackage{array}
\usepackage{multicol}
\usepackage{xcolor}
\usepackage{ragged2e}

% =========================
% TEMA Y COLORES
% =========================
\usetheme{Madrid}
\usecolortheme{default}
\setbeamertemplate{navigation symbols}{}

\definecolor{USTBlue}{HTML}{003B73}
\definecolor{USTLightBlue}{HTML}{DCEAF7}
\definecolor{USTGreen}{HTML}{4F772D}
\definecolor{USTGold}{HTML}{D9A404}
\definecolor{USTGray}{HTML}{4D4D4D}
\definecolor{USTSoft}{HTML}{F7F9FB}

\setbeamercolor{structure}{fg=USTBlue}
\setbeamercolor{title}{fg=white,bg=USTBlue}
\setbeamercolor{frametitle}{fg=white,bg=USTBlue}
\setbeamercolor{block title}{fg=white,bg=USTBlue}
\setbeamercolor{block body}{fg=black,bg=USTSoft}
\setbeamercolor{block title alerted}{fg=white,bg=USTGreen}
\setbeamercolor{block body alerted}{fg=black,bg=green!8}
\setbeamercolor{block title example}{fg=white,bg=USTGold}
\setbeamercolor{block body example}{fg=black,bg=yellow!10}
\setbeamercolor{item projected}{fg=white,bg=USTBlue}

% =========================
% PIE DE PÁGINA
% =========================
\setbeamertemplate{footline}{
	\leavevmode
	\hbox{
		\begin{beamercolorbox}[wd=.78\paperwidth,ht=2.5ex,dp=1ex,leftskip=0.3cm]{author in head/foot}
			\color{USTGray}\scriptsize Fundamentos de Matemática Básica --- Unidad I --- Clase 5: Unidades de medida
		\end{beamercolorbox}
		\begin{beamercolorbox}[wd=.22\paperwidth,ht=2.5ex,dp=1ex,rightskip=0.3cm plus1fil]{date in head/foot}
			\color{USTGray}\scriptsize \insertframenumber/\inserttotalframenumber
		\end{beamercolorbox}
	}
	\vskip0pt
}

% =========================
% COMANDOS
% =========================
\newcommand{\idea}[1]{
	\begin{block}{Idea clave}
		#1
	\end{block}
}

\newcommand{\ejemplo}[2]{
	\begin{exampleblock}{Ejemplo}
		#1
		
		\vspace{0.5em}
		\textbf{Resolución:} #2
	\end{exampleblock}
}

\newcommand{\alertacaja}[1]{
	\begin{alertblock}{Atención}
		#1
	\end{alertblock}
}

% =========================
% DATOS
% =========================
\title[MAT Básica]{Fundamentos de Matemática Básica}
\subtitle{Unidad I --- Clase 5: Unidades de medida}
\author{Profesor Luis Tulio González Alcaino}
\institute{Universidad Santo Tomás \\ Departamento Ciencias Básicas --- Talca}
\date{Semestre I -- 2026}

% =========================
% DOCUMENTO
% =========================
\begin{document}
	
	% -------------------------
	\begin{frame}
		\titlepage
	\end{frame}
	
	% -------------------------
	\begin{frame}{Objetivos de la clase}
		\begin{itemize}[<+->]
			\item Identificar el concepto de \textbf{magnitud} y su relación con la medición.
			\item Reconocer y convertir unidades de \textbf{longitud, masa, área, volumen y tiempo}.
			\item Resolver situaciones contextualizadas de conversión de unidades en \textbf{agronomía}.
			\item Detectar \textbf{patrones} e \textbf{inconsistencias} en cálculos con unidades.
		\end{itemize}
	\end{frame}
	
	% -------------------------
	\begin{frame}{Motivación}
		\begin{block}{Situación inicial}
			En agronomía, medir correctamente permite:
			\begin{itemize}[<+->]
				\item dosificar fertilizantes,
				\item calcular agua de riego,
				\item estimar superficies cultivadas,
				\item planificar tiempos de aplicación y cosecha.
			\end{itemize}
		\end{block}
		
		\pause
		\ejemplo
		{No basta decir: ``hay mucha agua en el estanque''.}
		{Debemos expresar una cantidad con unidad, por ejemplo:
			\[
			1500\ \text{L de agua}.
			\]
		}
	\end{frame}
	
	% -------------------------
	\begin{frame}{¿Qué es una magnitud?}
		\idea{Toda propiedad que puede medirse se denomina \textbf{magnitud}.}
		
		\begin{columns}[T]
			\column{0.48\textwidth}
			\begin{block}{Sí son magnitudes}
				\begin{itemize}[<+->]
					\item Longitud
					\item Masa
					\item Tiempo
					\item Área
					\item Volumen
				\end{itemize}
			\end{block}
			
			\column{0.48\textwidth}
			\begin{alertblock}{No son magnitudes}
				\begin{itemize}[<+->]
					\item Belleza del paisaje
					\item Sabor de un fruto
					\item Preferencia personal
				\end{itemize}
			\end{alertblock}
		\end{columns}
		
		\pause
		\begin{block}{Importancia en agronomía}
			Trabajar con magnitudes correctas permite dosificar, planificar y evitar errores productivos.
		\end{block}
	\end{frame}
	
	% -------------------------
	\begin{frame}{Clasificación de magnitudes}
		\begin{columns}[T]
			\column{0.48\textwidth}
			\begin{block}{Magnitudes fundamentales}
				Son aquellas que se miden directamente.
				\begin{itemize}[<+->]
					\item Longitud
					\item Masa
					\item Tiempo
				\end{itemize}
			\end{block}
			
			\column{0.48\textwidth}
			\begin{block}{Magnitudes derivadas}
				Se obtienen a partir de otras magnitudes.
				\begin{itemize}[<+->]
					\item Área
					\item Volumen
					\item Densidad
					\item Velocidad de riego
				\end{itemize}
			\end{block}
		\end{columns}
	\end{frame}
	
	% -------------------------
	\begin{frame}{Magnitudes frecuentes en agronomía}
		\small
		\begin{tabular}{>{\raggedright\arraybackslash}p{2.2cm} >{\raggedright\arraybackslash}p{2.2cm} >{\raggedright\arraybackslash}p{6cm}}
			\toprule
			\textbf{Magnitud} & \textbf{Unidad habitual} & \textbf{Aplicación en agronomía} \\
			\midrule
			Longitud & m, cm, km & Distancia entre surcos o entre plantas \\
			\pause
			Masa & g, kg & Fertilizantes, semillas y frutos \\
			\pause
			Tiempo & s, min, h & Riego, turnos de trabajo y cosecha \\
			\pause
			Área & m$^2$, ha & Superficie cultivada \\
			\pause
			Volumen & mL, L, m$^3$ & Agua, mezclas y estanques \\
			\bottomrule
		\end{tabular}
	\end{frame}
	
	% -------------------------
	\begin{frame}{Unidades básicas de longitud}
		\idea{
			Escala métrica:
			\[
			\text{km}\quad \text{hm}\quad \text{dam}\quad \text{m}\quad \text{dm}\quad \text{cm}\quad \text{mm}
			\]
			Cada paso hacia la derecha multiplica por \(10\).  
			Cada paso hacia la izquierda divide por \(10\).
		}
		
		\pause
		\ejemplo
		{La distancia entre plantas de tomate es de \(0{,}4\ \text{m}\). Exprésela en centímetros.}
		{
			\[
			0{,}4\ \text{m} = 0{,}4 \times 100\ \text{cm} = 40\ \text{cm}
			\]
			Por lo tanto, la distancia entre plantas es de \(\boxed{40\ \text{cm}}\).
		}
	\end{frame}
	
	% -------------------------
	\begin{frame}{Ejercicio resuelto de longitud (develado)}
		\begin{block}{Situación agronómica}
			En un huerto frutal, la separación entre dos hileras es de \(2{,}5\ \text{m}\).  
			Exprese esta longitud en centímetros.
		\end{block}
		
		\pause
		\textbf{Paso 1:} Identificamos la conversión
		\[
		1\ \text{m} = 100\ \text{cm}
		\]
		
		\pause
		\textbf{Paso 2:} Multiplicamos
		\[
		2{,}5 \times 100 = 250
		\]
		
		\pause
		\textbf{Respuesta:}
		\[
		2{,}5\ \text{m} = \boxed{250\ \text{cm}}
		\]
	\end{frame}
	
	% -------------------------
	\begin{frame}{Unidades básicas de masa}
		\idea{
			Escala métrica:
			\[
			\text{kg}\quad \text{hg}\quad \text{dag}\quad \text{g}\quad \text{dg}\quad \text{cg}\quad \text{mg}
			\]
			Cada paso hacia la derecha multiplica por \(10\).  
			Cada paso hacia la izquierda divide por \(10\).
		}
		
		\pause
		\ejemplo
		{Si una planta recibe \(250\ \text{g}\) de fertilizante, exprese la dosis en kilogramos.}
		{
			\[
			250\ \text{g} = \frac{250}{1000}\ \text{kg} = 0{,}25\ \text{kg}
			\]
			Entonces, la dosis equivale a \(\boxed{0{,}25\ \text{kg}}\).
		}
	\end{frame}
	
	% -------------------------
	\begin{frame}{Ejercicio resuelto de masa (develado)}
		\begin{block}{Situación agronómica}
			Un saco contiene \(4{,}2\ \text{kg}\) de semilla de trigo.  
			Exprese esta masa en gramos.
		\end{block}
		
		\pause
		\textbf{Paso 1:} Usamos la equivalencia
		\[
		1\ \text{kg} = 1000\ \text{g}
		\]
		
		\pause
		\textbf{Paso 2:} Convertimos
		\[
		4{,}2 \times 1000 = 4200
		\]
		
		\pause
		\textbf{Respuesta:}
		\[
		4{,}2\ \text{kg} = \boxed{4200\ \text{g}}
		\]
	\end{frame}
	
	% -------------------------
	\begin{frame}{Unidades de área}
		\idea{
			En área, cada salto entre unidades equivale a multiplicar o dividir por \(100\).
			\[
			\text{km}^2 \quad \text{hm}^2 \quad \text{dam}^2 \quad \text{m}^2 \quad \text{dm}^2 \quad \text{cm}^2 \quad \text{mm}^2
			\]
			Equivalencia importante:
			\[
			1\ \text{ha} = 10\,000\ \text{m}^2
			\]
		}
		
		\pause
		\ejemplo
		{Un predio tiene \(3{,}5\) hectáreas. Exprese esta superficie en metros cuadrados.}
		{
			\[
			3{,}5 \times 10\,000 = 35\,000
			\]
			Por lo tanto:
			\[
			3{,}5\ \text{ha} = \boxed{35\,000\ \text{m}^2}
			\]
		}
	\end{frame}
	
	% -------------------------
	\begin{frame}{Ejercicio resuelto de área (develado)}
		\begin{block}{Situación agronómica}
			Un terreno experimental de maíz tiene \(1{,}8\ \text{ha}\).  
			Exprese esta superficie en metros cuadrados.
		\end{block}
		
		\pause
		\textbf{Paso 1:} Recordamos que
		\[
		1\ \text{ha} = 10\,000\ \text{m}^2
		\]
		
		\pause
		\textbf{Paso 2:} Multiplicamos
		\[
		1{,}8 \times 10\,000 = 18\,000
		\]
		
		\pause
		\textbf{Respuesta:}
		\[
		1{,}8\ \text{ha} = \boxed{18\,000\ \text{m}^2}
		\]
	\end{frame}
	
	% -------------------------
	\begin{frame}{Unidades de volumen}
		\idea{
			En volumen, cada salto entre unidades equivale a multiplicar o dividir por \(1000\).
			\[
			\text{km}^3 \quad \text{hm}^3 \quad \text{dam}^3 \quad \text{m}^3 \quad \text{dm}^3 \quad \text{cm}^3 \quad \text{mm}^3
			\]
			Equivalencias útiles:
			\[
			1\ \text{L} = 1000\ \text{mL}
			\qquad
			1\ \text{L} = 1\ \text{dm}^3
			\]
		}
		
		\pause
		\ejemplo
		{Para una mezcla foliar se necesitan \(2{,}5\ \text{L}\) de agua. Exprese ese volumen en mililitros.}
		{
			\[
			2{,}5 \times 1000 = 2500
			\]
			Luego:
			\[
			2{,}5\ \text{L} = \boxed{2500\ \text{mL}}
			\]
		}
	\end{frame}
	
	% -------------------------
	\begin{frame}{Ejercicio resuelto de volumen (develado)}
		\begin{block}{Situación agronómica}
			En un invernadero se suministran \(3{,}5\ \text{L}\) de agua a una bandeja de almácigos.  
			Exprese esta cantidad en mililitros.
		\end{block}
		
		\pause
		\textbf{Paso 1:} Usamos la equivalencia
		\[
		1\ \text{L} = 1000\ \text{mL}
		\]
		
		\pause
		\textbf{Paso 2:} Convertimos
		\[
		3{,}5 \times 1000 = 3500
		\]
		
		\pause
		\textbf{Respuesta:}
		\[
		3{,}5\ \text{L} = \boxed{3500\ \text{mL}}
		\]
	\end{frame}
	
	% -------------------------
	\begin{frame}{Unidades de tiempo}
		\idea{
			Equivalencias:
			\[
			1\ \text{h} = 60\ \text{min}
			\qquad
			1\ \text{min} = 60\ \text{s}
			\]
			A diferencia de la longitud y la masa, el tiempo no cambia por potencias de \(10\).
		}
		
		\pause
		\ejemplo
		{Si un sistema de riego opera durante \(2{,}5\) horas, exprese ese tiempo en minutos.}
		{
			\[
			2{,}5 \times 60 = 150
			\]
			Entonces:
			\[
			2{,}5\ \text{h} = \boxed{150\ \text{min}}
			\]
		}
	\end{frame}
	
	% -------------------------
	\begin{frame}{Ejercicio resuelto de tiempo (develado)}
		\begin{block}{Situación agronómica}
			Un sistema de riego funciona durante \(2\ \text{h}\ 30\ \text{min}\).  
			Exprese el tiempo total en minutos.
		\end{block}
		
		\pause
		\textbf{Paso 1:} Convertimos las horas a minutos
		\[
		2\ \text{h} = 2 \times 60 = 120\ \text{min}
		\]
		
		\pause
		\textbf{Paso 2:} Sumamos los minutos restantes
		\[
		120 + 30 = 150
		\]
		
		\pause
		\textbf{Respuesta:}
		\[
		2\ \text{h}\ 30\ \text{min} = \boxed{150\ \text{min}}
		\]
	\end{frame}
	
	% -------------------------
	\begin{frame}{Operatoria con conversión de unidades}
		\begin{block}{Procedimiento general}
			\begin{enumerate}[<+->]
				\item Identificar la magnitud involucrada.
				\item Convertir todas las cantidades a una misma unidad.
				\item Realizar la operación.
				\item Expresar el resultado final con sentido agronómico.
			\end{enumerate}
		\end{block}
	\end{frame}
	
	% -------------------------
	\begin{frame}{Ejemplo 1: masa total de fertilizante}
		\begin{block}{Situación}
			Se aplican \(300\ \text{g}\) de fertilizante a cada planta en un sector con \(80\) plantas.
		\end{block}
		
		\pause
		\textbf{Paso 1:} Calculamos la masa total en gramos
		\[
		300 \times 80 = 24\,000\ \text{g}
		\]
		
		\pause
		\textbf{Paso 2:} Convertimos a kilogramos
		\[
		24\,000\ \text{g} = 24\ \text{kg}
		\]
		
		\pause
		\textbf{Resultado:}
		\[
		\boxed{24\ \text{kg}}
		\]
		Se requieren \(24\ \text{kg}\) de fertilizante.
	\end{frame}
	
	% -------------------------
	\begin{frame}{Ejemplo 2: área de cultivo}
		\begin{block}{Situación}
			Un predio tiene \(3{,}5\) hectáreas. Exprese esta superficie en metros cuadrados.
		\end{block}
		
		\pause
		\[
		1\ \text{ha} = 10\,000\ \text{m}^2
		\]
		
		\pause
		\[
		3{,}5 \times 10\,000 = 35\,000
		\]
		
		\pause
		\textbf{Resultado:}
		\[
		\boxed{35\,000\ \text{m}^2}
		\]
	\end{frame}
	
	% -------------------------
	\begin{frame}{Ejemplo 3: volumen total de agua}
		\begin{block}{Situación}
			Para una mezcla se necesitan \(4\) litros de agua por bandeja.  
			Si se preparan \(6\) bandejas, ¿qué volumen total se requiere?
		\end{block}
		
		\pause
		\textbf{Paso 1:} Multiplicamos
		\[
		4 \times 6 = 24\ \text{L}
		\]
		
		\pause
		\textbf{Paso 2:} Convertimos a mililitros
		\[
		24\ \text{L} = 24\,000\ \text{mL}
		\]
		
		\pause
		\textbf{Resultado:}
		\[
		\boxed{24\ \text{L}} \qquad \text{o} \qquad \boxed{24\,000\ \text{mL}}
		\]
	\end{frame}
	
	% -------------------------
	\begin{frame}{Ejemplo 4: tiempo de riego}
		\begin{block}{Situación}
			Una parcela se riega durante \(1\ \text{h}\ 45\ \text{min}\). Exprese este tiempo en minutos.
		\end{block}
		
		\pause
		\[
		1\ \text{h} = 60\ \text{min}
		\]
		
		\pause
		\[
		60 + 45 = 105\ \text{min}
		\]
		
		\pause
		\textbf{Resultado:}
		\[
		\boxed{105\ \text{min}}
		\]
	\end{frame}
	
	% -------------------------
	\begin{frame}{Patrones y relaciones en agronomía}
		\begin{block}{Patrones observables}
			\begin{itemize}[<+->]
				\item Si aumenta el número de plantas, aumenta la cantidad total de fertilizante.
				\item Si aumenta la superficie, aumenta el volumen de agua de riego.
				\item Si aumenta el tiempo de funcionamiento del sistema, aumenta el agua aplicada.
			\end{itemize}
		\end{block}
		
		\pause
		\idea{
			Relación general:
			\[
			\text{Cantidad total} = \text{dosis por unidad} \times \text{número de unidades}
			\]
		}
	\end{frame}
	
	% -------------------------
	\begin{frame}{Ejercicio resuelto de patrón (develado)}
		\begin{block}{Situación agronómica}
			Si cada planta de lechuga recibe \(150\ \text{mL}\) de agua y hay \(40\) plantas, ¿cuánta agua se necesita en total?
		\end{block}
		
		\pause
		\textbf{Paso 1:} Aplicamos la relación general
		\[
		\text{Cantidad total} = \text{dosis por planta} \times \text{número de plantas}
		\]
		
		\pause
		\[
		150 \times 40 = 6000\ \text{mL}
		\]
		
		\pause
		\textbf{Paso 2:} Convertimos a litros
		\[
		6000\ \text{mL} = 6\ \text{L}
		\]
		
		\pause
		\textbf{Respuesta:}
		\[
		\boxed{6000\ \text{mL} = 6\ \text{L}}
		\]
	\end{frame}
	
	% -------------------------
	\begin{frame}{Detección de inconsistencias}
		\begin{alertblock}{Errores frecuentes}
			\begin{itemize}[<+->]
				\item Sumar cantidades con distintas unidades sin convertir.
				\item Afirmar que \(1\ \text{ha} = 1000\ \text{m}^2\).
				\item Confundir litros con mililitros.
				\item No ajustar los insumos al aumentar la superficie cultivada.
			\end{itemize}
		\end{alertblock}
	\end{frame}
	
	% -------------------------
	\begin{frame}{Análisis de una inconsistencia}
		\begin{block}{Afirmación incorrecta}
			\[
			2\ \text{kg} + 500\ \text{g} = 2{,}500
			\]
		\end{block}
		
		\pause
		\textbf{Corrección:} primero convertimos a la misma unidad.
		\[
		2\ \text{kg} = 2000\ \text{g}
		\]
		
		\pause
		\[
		2000\ \text{g} + 500\ \text{g} = 2500\ \text{g}
		\]
		
		\pause
		\[
		2500\ \text{g} = 2{,}5\ \text{kg}
		\]
		
		\pause
		\alertacaja{La operación numérica solo tiene sentido si las unidades también son correctas.}
	\end{frame}
	
	% -------------------------
	\begin{frame}{Ejercicios propuestos}
		\small
		\begin{enumerate}[<+->]
			\item Convertir \(4{,}2\ \text{kg}\) de semilla a gramos.
			\item Convertir \(1{,}8\ \text{ha}\) a metros cuadrados.
			\item Convertir \(3{,}5\ \text{L}\) a mililitros.
			\item Expresar \(2\ \text{h}\ 30\ \text{min}\) en minutos.
			\item Detectar el error en la afirmación:
			\[
			1\ \text{ha} = 1000\ \text{m}^2
			\]
		\end{enumerate}
	\end{frame}
	
	% -------------------------
	\begin{frame}{Ejercicios resueltos}
		\small
		\begin{enumerate}[<+->]
			\item \[
			4{,}2\ \text{kg} = 4200\ \text{g}
			\]
			
			\item \[
			1{,}8\ \text{ha} = 18\,000\ \text{m}^2
			\]
			
			\item \[
			3{,}5\ \text{L} = 3500\ \text{mL}
			\]
			
			\item \[
			2\ \text{h}\ 30\ \text{min} = 150\ \text{min}
			\]
			
			\item La afirmación es incorrecta porque:
			\[
			1\ \text{ha} = 10\,000\ \text{m}^2
			\]
		\end{enumerate}
	\end{frame}
	
	% -------------------------
	\begin{frame}{Ejercicios contextualizados a agronomía}
		\small
		\begin{enumerate}[<+->]
			\item En un cultivo experimental se debe aplicar un fertilizante en una dosis de \(500\ \text{mg}\) por planta. ¿Cuántos gramos equivalen a esta dosis?
			\item En un invernadero se recomienda suministrar \(2\ \text{L}\) de agua al día a una bandeja de almácigos. ¿Cuántos mililitros equivalen a esta cantidad?
			\item Para el tratamiento de semillas, se requiere aplicar \(750\ \mu\text{g}\) de un regulador de crecimiento por muestra. ¿Cuántos miligramos equivalen a esta dosis?
			\item En un sistema de riego automatizado, una parcela recibe agua tres veces al día, con intervalos de 6 horas. Si el primer riego comienza a las 8:00, ¿a qué hora deben realizarse los otros dos riegos?
			\item En un sistema de fertirrigación, una solución nutritiva se aplica a razón de \(100\ \text{mL}\) por hora. ¿Cuántos litros se aplican en 8 horas?
		\end{enumerate}
	\end{frame}
	
	% -------------------------
	\begin{frame}{Soluciones contextualizadas (1 y 2)}
		\small
		\begin{block}{1. Fertilizante por planta}
			Se sabe que:
			\[
			1\ \text{g} = 1000\ \text{mg}
			\]
			\pause
			Entonces:
			\[
			500\ \text{mg}=\frac{500}{1000}\ \text{g}=0{,}5\ \text{g}
			\]
			\pause
			\textbf{Respuesta:}
			Cada planta requiere \(\boxed{0{,}5\ \text{g}}\) de fertilizante.
		\end{block}
		
		\pause
		\begin{block}{2. Riego en almácigos}
			Se sabe que:
			\[
			1\ \text{L} = 1000\ \text{mL}
			\]
			\pause
			Entonces:
			\[
			2\ \text{L}=2\times 1000=2000\ \text{mL}
			\]
			\pause
			\textbf{Respuesta:}
			Cada bandeja recibe \(\boxed{2000\ \text{mL}}\) de agua diaria.
		\end{block}
	\end{frame}
	
	% -------------------------
	\begin{frame}{Soluciones contextualizadas (3 y 4)}
		\small
		\begin{block}{3. Regulador de crecimiento}
			Se sabe que:
			\[
			1\ \text{mg} = 1000\ \mu\text{g}
			\]
			\pause
			Entonces:
			\[
			750\ \mu\text{g}=\frac{750}{1000}\ \text{mg}=0{,}75\ \text{mg}
			\]
			\pause
			\textbf{Respuesta:}
			La dosis aplicada es de \(\boxed{0{,}75\ \text{mg}}\) por muestra.
		\end{block}
		
		\pause
		\begin{block}{4. Sistema de riego}
			Primer riego:
			\[
			8{:}00
			\]
			\pause
			Segundo riego:
			\[
			8{:}00 + 6\ \text{h} = 14{:}00
			\]
			\pause
			Tercer riego:
			\[
			14{:}00 + 6\ \text{h} = 20{:}00
			\]
			\pause
			\textbf{Respuesta:}
			Los otros dos riegos deben realizarse a las \(\boxed{14{:}00}\) y \(\boxed{20{:}00}\).
		\end{block}
	\end{frame}
	
	% -------------------------
	\begin{frame}{Solución contextualizada (5)}
		\small
		\begin{block}{5. Fertirrigación}
			La solución se aplica a razón de:
			\[
			100\ \text{mL/h}
			\]
			Durante \(8\) horas:
			\[
			100 \times 8 = 800\ \text{mL}
			\]
			\pause
			Convertimos a litros:
			\[
			800\ \text{mL} = 0{,}8\ \text{L}
			\]
			\pause
			\textbf{Respuesta:}
			Se aplican \(\boxed{0{,}8\ \text{L}}\) de solución nutritiva en la jornada.
		\end{block}
	\end{frame}
	
	% -------------------------
	\begin{frame}{Síntesis}
		\begin{itemize}[<+->]
			\item Medir es comparar una cantidad con una unidad.
			\item En agronomía, masa, longitud, tiempo, área y volumen aparecen constantemente.
			\item Convertir correctamente las unidades permite tomar decisiones adecuadas.
			\item Detectar patrones ayuda a anticipar resultados.
			\item Detectar inconsistencias evita errores técnicos y económicos.
		\end{itemize}
	\end{frame}
	
	% -------------------------
	\begin{frame}{Cierre}
		\centering
		\vspace{1cm}
		{\Huge Fin de la clase}
		
		\vspace{0.6cm}
		{\Large ¿Preguntas?}
	\end{frame}
	
\end{document}